\section{Related Work}
\label{sec:related_work}

The computational bottleneck of Transformer models~\cite{vaswani2017attention} originates from two distinct dimensions: the quadratic memory and time complexity $\mathcal{O}(N^2)$ of full self-attention with respect to sequence length $N$, and the cubic parameter and compute footprint $\mathcal{O}(d^2)$ of projection matrices within multi-head attention and feedforward sublayers. Prior research has approached these dual challenges along separate trajectories.

\subsection{Low-Rank Matrix Factorization and Post-Training SVD}
Low-rank matrix decomposition derives from the classical Eckart-Young-Mirsky theorem~\cite{eckart1936approximation}, which proves that truncated Singular Value Decomposition (SVD) provides the optimal rank-$k$ approximation under the Frobenius and spectral norms. In deep language modeling, randomized SVD algorithms~\cite{halko2011structure} have been adopted to compress large weight matrices without re-training from scratch. Recent post-training compression techniques such as ASVD~\cite{yuan2023asvd} and SVD-LLM~\cite{wang2024svdllm} incorporate activation awareness and truncation-aware whitening, mitigating severe accuracy degradation caused by non-uniform feature channel magnitudes. Similarly, Kronecker decomposition~\cite{edalati2022kronecker} factorizes weight projections into tensor products to compress pre-trained representations. However, unconstrained dynamic rank selection often induces numerical singularities or under-allocates capacity to critical classification heads, leading to representation collapse on fine-grained discrimination tasks. Furthermore, weight matrix factorization alone leaves the quadratic $\mathcal{O}(N^2)$ sequence scaling of self-attention unmitigated.

\subsection{Sub-Quadratic, Linear, and Kernelized Attention}
To overcome the $\mathcal{O}(N^2)$ sequence bottleneck, extensive literature explores linear-time attention approximations. Linear Transformers~\cite{katharopoulos2020transformers} decompose the softmax attention kernel using right-associative feature maps $\phi(Q)\phi(K)^T V$, transforming attention computation into an autoregressive recurrent accumulation with $\mathcal{O}(N)$ computational complexity and $\mathcal{O}(1)$ step memory. Performers~\cite{choromanski2021rethinking} construct positive orthogonal random feature approximations (FAVOR+) to provably track softmax distributions without kernel divergence. Linformer~\cite{wang2020linformer} projects sequence length dimensions into low-dimensional subspaces via learned projection matrices, while Nystr\"omformer~\cite{xiong2021nystromformer} adopts landmark segment clustering coupled with iterative Moore-Penrose pseudoinverse approximations. cosFormer~\cite{qin2022cosformer} incorporates cosine reweighting to reinforce locality while maintaining linear time complexity. Despite their algorithmic efficiency, kernel-linearized models often experience entropy diffusion on cross-sentence alignment benchmarks (e.g., Natural Language Inference) due to the absence of the exponential softmax sharpeness, and standard implementations retain full-rank parameter projections.

\subsection{Structured State Space Models and Linear Recurrence}
State Space Models (SSMs) provide an alternative framework for long-sequence modeling. S4~\cite{gu2022efficiently} parameterizes continuous-time state representations via the High-Order Polynomial Projection Operators (HiPPO) framework, capturing long-range dependencies in $\mathcal{O}(N \log N)$ compute. Mamba~\cite{gu2023mamba} introduces data-dependent selective state spaces, varying discretization timescales $\Delta$ dynamically across tokens to match the in-context filtering capabilities of attention while executing in linear time $\mathcal{O}(N)$ with bounded state footprint. Recent extensions such as HGRN-2~\cite{qin2024hgrn2} expand state representations using outer-product recurrence. However, as demonstrated empirically and theoretically by Vision Mamba~\cite{zhu2024visionmamba}, causal 1D state-space scans possess an inherent unidirectional inductive bias: when substituted naively into bidirectional masked language encoders (e.g., DistilBERT~\cite{sanh2019distilbert}), the lack of future-context conditioning leads to degenerate representation collapse unless dual-pass scanning or bidirectional pre-training is implemented.

\subsection{Mixture-of-Experts and High-Capacity Hybrid Systems}
Rather than compressing physical parameters, conditional computation techniques expand total capacity while bounding active FLOPs. Switch Transformers~\cite{fedus2022switch} and Mixtral of Experts~\cite{jiang2024mixtral} route tokens across multiple feedforward subnetworks using top-$k$ gating. Expert Choice MoE~\cite{zhou2022expertchoice} inverts this routing topology to achieve deterministic load balancing without dropped tokens. Hybrid architectures interleave self-attention with recurrent state spaces and MoE blocks: Jamba~\cite{lieber2024jamba} combines 1:7 Transformer-to-Mamba ratios with feedforward experts, Samba~\cite{ren2024samba} merges selective SSMs with sliding window attention, and Hymba~\cite{dong2024hymba} incorporates parallel attention and state heads. As detailed in our Appendix, while these hybrid architectures exhibit high expressive capacity, their physical parameter footprints expand by up to $+170.25\%$, requiring substantial VRAM allocations ($>3.8$ GB) that diverge from the constraints of edge-deployed parameter compression.

\subsection{Parameter-Efficient Fine-Tuning and Pruning}
Post-hoc adaptation of compressed representations relies extensively on parameter-efficient tuning frameworks. LoRA~\cite{hu2022lora} and QLoRA~\cite{dettmers2023qlora} inject trainable rank-$r$ residual adapters into frozen linear projections, drastically reducing optimization memory. In post-training pruning, one-shot methods such as SparseGPT~\cite{frantar2023sparsegpt}, Wanda~\cite{sun2023wanda}, and SliceGPT~\cite{ashkboos2024slicegpt} remove redundant matrix channels based on second-order Hessian approximations or activation magnitudes. In serving environments, PagedAttention~\cite{kwon2023pagedattention} and StreamingLLM~\cite{xiao2024streamingllm} mitigate memory fragmentation and sustain infinite-length generation via attention sinks. Post-training quantization frameworks like SmoothQuant~\cite{xiao2023smoothquant}, AWQ~\cite{lin2024awq}, and GPTQ~\cite{frantar2023gptq} map FP16 tensors into INT8/INT4 domains by balancing activation outliers against weight scales.

\subsection{Positioning and Synthesis of Proposed Framework}
Our work synthesizes these previously isolated domains into a unified framework: \textbf{SVD Linear Attention}. Whereas prior SVD methods solely reduce weight parameters while retaining quadratic $\mathcal{O}(N^2)$ memory, and linear attention methods solely reduce token complexity while retaining dense projection matrices, our framework couples exact truncated SVD projection factorization with recurrent state-space linear attention. Furthermore, by introducing the Adaptive Energy-Guided Iterative SVD (AEGIS) formulation with rank floor guarantees ($r_{min} \ge 16$) and hardware-stabilized high-order Schulz polynomial matrix inversions, our approach eliminates both the numerical instability of iterative pseudoinverses in half-precision arithmetic and the rank starvation of classification heads. Across 217 rigorously evaluated attention and compression modules, this balanced design establishes a Pareto-optimal operating point, retaining over $91.0\%$ accuracy while providing physical parameter compression ($1.15\times$ to $1.25\times$) and strictly bounded memory scaling.
